\documentclass[12pt,a4paper]{article}

\usepackage[utf8]{vietnam}
\usepackage[utf8]{inputenc}
\usepackage[left=2.5cm,right=2.5cm,top=2.5cm,bottom=2.5cm]{geometry}

\usepackage{amsmath,amssymb,amsfonts,amsthm,mathtools}
\usepackage{stmaryrd}
\usepackage[hidelinks]{hyperref}

\theoremstyle{definition}
\newtheorem{dinhnghia}{Định nghĩa}
\newtheorem*{ghichu}{Ghi chú}

\newcommand{\plays}{\mathrel{\mathsf{plays}}}
\newcommand{\playsd}{\mathrel{\mathsf{plays}_{d}}}

\title{Cú pháp và Ngữ nghĩa Hình thức của Object Model Mở rộng}
\date{}

\begin{document}
\maketitle


% ============================================================
\section{Cú pháp hình thức của Object Model mở rộng}
% ============================================================

Giả sử
\[
N\subseteq A^{+},
\qquad
\Sigma=(T,\Omega).
\]

$N$ là tập tên và $\Sigma$ là signature của các kiểu dữ liệu.

\begin{ghichu}
Tài liệu này là một mở rộng của Object Model nền trong
\texttt{UML\_formal.tex}. Mọi định nghĩa và điều kiện well-formedness của
mô hình nền vẫn được giữ nguyên nếu không được định nghĩa lại ở đây; một điều
kiện mang cùng nhãn trong tài liệu này thay thế điều kiện tương ứng của mô hình
nền. Vì vậy, điều kiện cấm chu trình composition giữa các Group (WF-8 của mô
hình nền) tiếp tục có hiệu lực mà không cần được lặp lại.
\end{ghichu}


% ============================================================
\subsection{Class}
% ============================================================

\[
Class = Entity\cup Role\cup Group\subseteq N,
\qquad
|Class|<\infty.
\]


\[
Entity\cap Role
=
Entity\cap Group
=
Role\cap Group
=
\varnothing.
\]


\[
\forall c\in Class,\qquad
\exists!\,t_c\in T.
\]

Mỗi lớp $c$ xác định duy nhất một kiểu đối tượng $t_c$.


% ============================================================
\subsection{Attribute}
% ============================================================

\[
Att_c
\subseteq
\{\,a:t_c\to t\mid a\in N,\ t\in T\,\},
\qquad c\in Class.
\]


\[
\forall t,t'\in T:
\quad
\bigl(
a:t_c\to t\in Att_c
\land
a:t_c\to t'\in Att_c
\bigr)
\Rightarrow
t=t'.
\]

Một tên attribute trong cùng một lớp chỉ xác định một kiểu.


% ============================================================
\subsection{Association}
% ============================================================

\subsubsection{Association}

\[
Assoc\subseteq N,
\qquad
|Assoc|<\infty.
\]

$Assoc$ là tập hữu hạn các association.

\[
associates:
\left\{
\begin{array}{rcl}
Assoc & \to & Class^{+}\\
as & \mapsto & \langle c_1,\ldots,c_n\rangle,\quad n\geq2
\end{array}
\right.
\]

$associates(as)$ xác định các lớp tham gia association $as$.


% ------------------------------------------------------------
\subsubsection{Role name}
% ------------------------------------------------------------

\[
rolenames:
\left\{
\begin{array}{rcl}
Assoc & \to & N^{+}\\
as & \mapsto & \langle \rho_1,\ldots,\rho_n\rangle
\end{array}
\right.
\]

$rolenames(as)$ xác định role name tại các đầu của association $as$.

\[
associates(as)=\langle c_1,\ldots,c_n\rangle
\Rightarrow
|rolenames(as)|=n.
\]

Số role name bằng số đầu của association.

\[
\forall i,j\in\{1,\ldots,n\}:
\qquad
i\neq j\Rightarrow \rho_i\neq \rho_j.
\]

Các role name trong cùng một association phải khác nhau.


% ------------------------------------------------------------
\subsubsection{Participating association}
% ------------------------------------------------------------

\[
participating:
\left\{
\begin{array}{rcl}
Class & \to & \mathcal{P}(Assoc)\\
c & \mapsto &
\{\,as\mid
as\in Assoc
\land
associates(as)=\langle c_1,\ldots,c_n\rangle
\\
&&
\qquad\land
\exists i\in\{1,\ldots,n\}:c_i=c
\,\}
\end{array}
\right.
\]

$participating(c)$ là tập các association mà lớp $c$ tham gia.


% ------------------------------------------------------------
\subsubsection{Navigable ends}
% ------------------------------------------------------------

\[
navends:
\left\{
\begin{array}{rcl}
Class\times Assoc & \to & \mathcal{P}(N)\\
(c,as) & \mapsto &
\{\,\rho\mid
associates(as)=\langle c_1,\ldots,c_n\rangle
\\
&&
\qquad\land
rolenames(as)=\langle \rho_1,\ldots,\rho_n\rangle
\\
&&
\qquad\land
\exists i,j\in\{1,\ldots,n\}:
(i\neq j\land c_i=c\land \rho_j=\rho)
\,\}
\end{array}
\right.
\]

$navends(c,as)$ là tập các role name có thể điều hướng tới từ $c$ qua $as$.

\[
navends:
\left\{
\begin{array}{rcl}
Class & \to & \mathcal{P}(N)\\
c & \mapsto &
\displaystyle
\bigcup_{as\in participating(c)}
navends(c,as)
\end{array}
\right.
\]

$navends(c)$ chứa toàn bộ các role name có thể điều hướng từ $c$.


% ------------------------------------------------------------
\subsubsection{Multiplicity}
% ------------------------------------------------------------

\[
multiplicities:
\left\{
\begin{array}{rcl}
Assoc
& \to &
\bigl(
\mathcal{P}(\mathbb{N}_0)\setminus\{\varnothing\}
\bigr)^{+}
\\
as
& \mapsto &
\langle M_1,\ldots,M_n\rangle
\end{array}
\right.
\]

$multiplicities(as)$ xác định multiplicity tại từng đầu association.

\[
\varnothing\neq M_i\subseteq\mathbb{N}_0,
\qquad
M_i\neq\{0\},
\qquad
1\leq i\leq n.
\]

Mỗi multiplicity là một tập cardinality hợp lệ và khác $\{0\}$.

\[
|multiplicities(as)|
=
|rolenames(as)|
=
|associates(as)|.
\]

Mỗi association end có đúng một role name và một multiplicity.


% ============================================================


% ============================================================
\subsection{Aggregation and Composition}
% ============================================================

\[
Aggregation\subseteq Assoc,
\qquad
Composition\subseteq Assoc.
\]

$Assoc$ là khái niệm cha; $Aggregation$ và $Composition$ là hai loại
association chuyên biệt và sử dụng cùng cấu trúc end, role name và
multiplicity của association.

$Composition$ là tập các association có ngữ nghĩa composition.

\[
\forall as\in Assoc:
\]
\[
\left(
\begin{array}{c}
associates(as)=\langle c_1,\ldots,c_n\rangle
\\[1mm]
\land
\exists i\neq j:\ c_i,c_j\in Role\cup Group
\end{array}
\right)
\iff
\left(
\begin{array}{c}
n=2
\\
\land\ as\in Composition
\\
\land\ c_1\in Group
\\
\land\ c_2\in Role\cup Group
\\
\land\ c_1\neq c_2
\end{array}
\right).
\tag{WF-6}
\]

Association giữa hai lớp thuộc $Role\cup Group$ chỉ có thể là composition từ $Group$ tới $Role$ hoặc $Group$.

Đặc biệt, một association Role--Role không thỏa WF-6: vế trái đúng nhưng vế
phải sai vì lớp đầu tiên không thuộc $Group$. Do đó profile ACL này không cho
phép association Role--Role.

\[
\forall x\in Role\cup Group:
\]
\[
\left|
\left\{
g\in Group
\;\middle|\;
\exists as\in Composition:\;
associates(as)=\langle g,x\rangle
\right\}
\right|
\leq1.
\tag{WF-7}
\]

Mỗi lớp thuộc $Role\cup Group$ có nhiều nhất một $Group$ sở hữu thông qua
Composition; không cần đưa thêm quan hệ $Owner$ vào metamodel.



% ============================================================
\subsection{Generalization}
% ============================================================


\[
\prec\ \subseteq
(Entity\times Entity)
\cup
(Group\times Group).
\tag{WF-9}
\]

Generalization chỉ được định nghĩa giữa các lớp cùng loại.

\[
parents:
\left\{
\begin{array}{rcl}
Class & \to & \mathcal{P}(Class)\\
c & \mapsto &
\{\,c'\mid c'\in Class\land c\prec c'\,\}
\end{array}
\right.
\]

$parents(c)$ chứa toàn bộ các lớp cha của $c$.

\[
Att_c^{*}
=
Att_c
\cup
\bigcup_{c'\in parents(c)}
Att_{c'},
\qquad
c\in Class.
\]

$Att_c^{*}$ chứa các attribute trực tiếp và các attribute kế thừa.

\[
navends^{*}(c)
=
navends(c)
\cup
\bigcup_{c'\in parents(c)}
navends(c').
\]

$navends^{*}(c)$ chứa các navigable end trực tiếp và kế thừa.

\[
FD_c
=
\bigl(
Att_c^{*},
navends^{*}(c)
\bigr).
\]

$FD_c$ là full descriptor của lớp $c$.



% ============================================================


% ============================================================
\subsection{Role Inheritance}
% ============================================================

Quan hệ kế thừa giữa các Role được xác định bởi

\[
\prec_r\ \subseteq Role\times Role.
\tag{WF-10}
\]

Quan hệ $\prec_r$ là một strict partial order:

\[
\forall r\in Role:
\qquad
\neg(r\prec_r r),
\]

\[
\forall r_1,r_2,r_3\in Role:
\quad
(r_1\prec_r r_2\land r_2\prec_r r_3)
\Rightarrow
r_1\prec_r r_3.
\]

\[
roleparents:
\left\{
\begin{array}{rcl}
Role & \to & \mathcal{P}(Role)\\
r & \mapsto &
\{\,r'\mid r'\in Role\land r\prec_r r'\,\}
\end{array}
\right.
\]

$roleparents(r)$ chứa toàn bộ các Role cha của $r$.

\[
Att_r^{*}
=
Att_r
\cup
\bigcup_{r'\in roleparents(r)}
Att_{r'},
\qquad
r\in Role.
\]

\[
navends^{*}(r)
=
navends(r)
\cup
\bigcup_{r'\in roleparents(r)}
navends(r'),
\qquad
r\in Role.
\]

\[
FD_r
=
\bigl(
Att_r^{*},
navends^{*}(r)
\bigr),
\qquad
r\in Role.
\]

\[
children_r:
\left\{
\begin{array}{rcl}
Role & \to & \mathcal{P}(Role)\\
r & \mapsto &
\{\,r'\in Role
\mid
r'\prec_r r
\land
\neg\exists r''\in Role:
(r'\prec_r r''\land r''\prec_r r)
\,\}
\end{array}
\right.
\]


% ============================================================
\subsection{Role Compatibility}
% ============================================================

\[
Compatibility\subseteq Role\times Role.
\]

Quan hệ compatibility giữa các Role là đối xứng:

\[
\forall r_1,r_2\in Role:
\qquad
(r_1,r_2)\in Compatibility
\Longleftrightarrow
(r_2,r_1)\in Compatibility.
\]

$Compatibility$ là quan hệ ở mức đặc tả. Bộ chuyển đổi dịch quan hệ này
thành các biểu thức OCL; nó không sinh thêm class, object hoặc association
link để biểu diễn compatibility trong object model đích.




\subsection{Well-Formedness}
% ============================================================

\paragraph{WF-1.}

\[
\forall c\in Class,\;
\forall as_1,as_2\in participating(c):
\]
\[
as_1\neq as_2
\Rightarrow
navends(c,as_1)\cap navends(c,as_2)
=
\varnothing.
\tag{WF-1}
\]

Các navigable role name từ cùng một lớp phải phân biệt giữa các association.


\paragraph{WF-2.}

\[
\forall c\in Class,\;
\forall
(a:t_{c_1}\to t_1),
(a':t_{c_2}\to t_2)
\in Att_c^{*}:
\]
\[
a=a'
\Rightarrow
(c_1=c_2\land t_1=t_2).
\tag{WF-2}
\]

Một attribute name chỉ được định nghĩa tại một lớp trong full descriptor.


% WF-3 của Object Model gốc liên quan đến operation và không được sử dụng ở đây.


\paragraph{WF-4.}

\[
\forall c\in Class,\;
\forall c_1,c_2\in parents(c)\cup\{c\}:
\]
\[
c_1\neq c_2
\Rightarrow
navends(c_1)\cap navends(c_2)
=
\varnothing.
\tag{WF-4}
\]

Một navigable role name chỉ được định nghĩa tại một lớp trong hierarchy.


\paragraph{WF-5.}

\[
\forall c\in Class,\;
\forall(a:t_{c'}\to t)\in Att_c^{*},\;
\forall \rho\in navends^{*}(c):
\]
\[
a\neq \rho.
\tag{WF-5}
\]

Attribute name và navigable role name trong cùng full descriptor không được trùng nhau.

\paragraph{WF-R2.}

\[
\forall r\in Role,\;
\forall
(a:t_{r_1}\to t_1),
(a':t_{r_2}\to t_2)
\in Att_r^{*}:
\]
\[
a=a'
\Rightarrow
(r_1=r_2\land t_1=t_2).
\tag{WF-R2}
\]

Một attribute name chỉ được định nghĩa tại một Role trong full descriptor.


\paragraph{WF-R4.}

\[
\forall r\in Role,\;
\forall r_1,r_2\in roleparents(r)\cup\{r\}:
\]
\[
r_1\neq r_2
\Rightarrow
navends(r_1)\cap navends(r_2)
=
\varnothing.
\tag{WF-R4}
\]

Một navigable role name chỉ được định nghĩa tại một Role trong hierarchy.


\paragraph{WF-R5.}

\[
\forall r\in Role,\;
\forall(a:t_{r'}\to t)\in Att_r^{*},\;
\forall \rho\in navends^{*}(r):
\]
\[
a\neq \rho.
\tag{WF-R5}
\]

Attribute name và navigable role name trong cùng full descriptor của Role
không được trùng nhau.



% ============================================================


\section{Interpretation of the Extended Object Model}
\subsection{Objects}
% ============================================================

\[
oid(c)=\{c_1,c_2,\ldots\},
\qquad
c\in Class.
\]

$oid(c)$ là tập vô hạn các định danh đối tượng tiềm năng của lớp $c$.

\[
I_{Class}(c)
=
\begin{cases}
\displaystyle
\bigcup
\left\{
oid(c')
\;\middle|\;
c'\in Entity\cup Group
\land
c'\preceq c
\right\},
&
c\in Entity\cup Group,
\\[5mm]
oid(c),
&
c\in Role.
\end{cases}
\]

$Entity$ và $Group$ sử dụng domain inclusion theo quan hệ
generalization $\prec$, trong khi mỗi $Role$ được diễn giải trên miền
object riêng.


% ============================================================
\subsection{Interpretation of Generalization}
% ============================================================

\[
\forall c_1,c_2\in Entity\cup Group:
\qquad
c_1\prec c_2
\Longrightarrow
I_{Class}(c_1)
\subseteq
I_{Class}(c_2).
\]

Quan hệ generalization $\prec$ của $Entity$ và $Group$ được diễn giải
bằng domain inclusion.


% ============================================================
\subsection{Interpretation of Role Inheritance}
% ============================================================

\[
\forall r\in Role:
\qquad
I_{Class}(r)=oid(r).
\]

Quan hệ Role inheritance $\prec_r$ không được diễn giải bằng domain
inclusion; do đó $r\prec_r r'$ không làm cho object của $r$ trở thành
object của $r'$. Quan hệ giữa các instance của các Role trong hierarchy
được biểu diễn riêng trong system state.


% ============================================================
\subsection{Links}
% ============================================================

\[
associates(as)=\langle c_1,\ldots,c_n\rangle
\Longrightarrow
I_{Assoc}(as)
=
I_{Class}(c_1)
\times\cdots\times
I_{Class}(c_n).
\]

$I_{Assoc}(as)$ là miền của tất cả các link khả dĩ của association $as$.

\[
l_{as}\in I_{Assoc}(as).
\]

Mỗi phần tử của $I_{Assoc}(as)$ là một link của association $as$.


% ============================================================


\subsection{System State}
% ============================================================

Một trạng thái hệ thống của model $M$ là

\[
\sigma(M)
=
\bigl(
\sigma_{Class},
\sigma_{Att},
\sigma_{Assoc},
\sigma_{Play}
\bigr).
\]

$\sigma(M)$ xác định object, attribute, association link và các
play-link giữa các Role instance tại một thời điểm.

$Compatibility$ không bổ sung thành phần vào $\sigma(M)$; các biểu thức
OCL được sinh từ $Compatibility$ được đánh giá trực tiếp trên trạng thái này.

\begin{enumerate}
\renewcommand{\labelenumi}{(\roman{enumi})}

% ------------------------------------------------------------
\item

\[
\forall c\in Class:
\qquad
\sigma_{Class}(c)\subset oid(c),
\qquad
|\sigma_{Class}(c)|<\infty.
\]

$\sigma_{Class}(c)$ chứa các object của lớp $c$ đang tồn tại trong state.


% ------------------------------------------------------------
\item

\[
\forall c\in Entity\cup Group,\;
\forall a:t_c\to t\in Att_c^{*}:
\]

\[
\sigma_{Att}(a):
\sigma_{Class}(c)\to I(t).
\]

Các object của $Entity$ và $Group$ sử dụng các attribute trực tiếp
và kế thừa theo ngữ nghĩa domain inclusion.

\[
\forall r\in Role,\;
\forall a:t_r\to t\in Att_r:
\]

\[
\sigma_{Att}(a):
\sigma_{Class}(r)\to I(t).
\]

Mỗi Role instance chỉ lưu các attribute được khai báo trực tiếp
tại chính Role đó.


% ------------------------------------------------------------
\item

\[
\forall as\in Assoc:
\qquad
\sigma_{Assoc}(as)
\subset
I_{Assoc}(as),
\qquad
|\sigma_{Assoc}(as)|<\infty.
\]

$\sigma_{Assoc}(as)$ chứa các link của association $as$ đang tồn tại.

\[
associates(as)=\langle c_1,\ldots,c_n\rangle
\]

\[
\Longrightarrow
\forall i\in\{1,\ldots,n\},\;
\forall l\in\sigma_{Assoc}(as):
\]

\[
\left|
\left\{
l'
\;\middle|\;
l'\in\sigma_{Assoc}(as)
\land
\bar{\pi}_i(l')
=
\bar{\pi}_i(l)
\right\}
\right|
\in
\pi_i\bigl(multiplicities(as)\bigr).
\]

Trong đó $\pi_i(l)$ là phép chiếu thành phần thứ $i$, còn
$\bar{\pi}_i(l)$ là phép chiếu tất cả các thành phần ngoại trừ
thành phần thứ $i$.


% ------------------------------------------------------------
\item

\[
\forall r\in Role,\;
\forall r'\in children_r(r):
\]

\[
\sigma_{Play}(r,r'):
\sigma_{Class}(r)
\to
\mathcal{P}\bigl(\sigma_{Class}(r')\bigr).
\]

\[
\forall y\in\sigma_{Class}(r'):
\qquad
\left|
\left\{
x\in\sigma_{Class}(r)
\;\middle|\;
y\in\sigma_{Play}(r,r')(x)
\right\}
\right|
=1.
\]



% ============================================================
\section{Translation of Role Property Access to Standard OCL}
% ============================================================

Ngữ pháp và ngữ nghĩa của OCL được giữ nguyên. Trước khi một biểu thức
được đưa sang OCL, bộ chuyển đổi làm tường minh việc truy cập các property
kế thừa của Role thông qua navigation \texttt{playOf}.

Với

\[
r\in children_r(r'),
\]

bộ chuyển đổi sinh một association phụ trợ $as^{play}_{r',r}$ sao cho

\[
associates(as^{play}_{r',r})
=
\langle r',r\rangle,
\]

\[
rolenames(as^{play}_{r',r})
=
\langle \mathtt{playOf},\rho^{fresh}_{r',r}\rangle,
\]

và các link của association này được xác định từ $\sigma_{Play}$ bởi

\[
\sigma'_{Assoc}(as^{play}_{r',r})
=
\left\{
(x,y)
\;\middle|\;
x\in\sigma_{Class}(r')
\land
y\in\sigma_{Play}(r',r)(x)
\right\}.
\]

Do đó \texttt{playOf} là một navigation OCL thông thường từ instance
của Role con $r$ tới instance tương ứng của Role cha $r'$.

Nếu property $p$ không được khai báo tại $r$ nhưng được khai báo tại
Role cha $r'$, thì bộ chuyển đổi thực hiện

\[
e.p
\rightsquigarrow
e.\mathtt{playOf}.p.
\]

Việc biến đổi được lặp lại khi property được khai báo ở Role cao hơn
trong hierarchy. Chẳng hạn,

\[
Director\in children_r(Employee),
\qquad
Employee\in children_r(Person),
\]

và $email$ được khai báo tại $Person$, thì

\[
d.email
\rightsquigarrow
d.\mathtt{playOf}.\mathtt{playOf}.email.
\]

Sau bước biến đổi, biểu thức thu được chỉ sử dụng các navigation và
property access thông thường và được đánh giá trực tiếp bằng OCL chuẩn.


\end{document}
